\documentclass[final]{beamer}

\usepackage[orientation=landscape,size=a0,scale=1.15]{beamerposter}
\usepackage{graphicx}
\usepackage{booktabs}
\usepackage{amsmath}
\usepackage{xcolor}
\usepackage{tikz}
\usepackage{array}
\usepackage{colortbl}

% ── CMU palette ──────────────────────────────────────────────────────────────
\definecolor{cmuRed}{HTML}{CC0000}
\definecolor{cmuGray}{HTML}{444444}
\definecolor{lightGray}{HTML}{F4F4F4}
\definecolor{accentBlue}{HTML}{2C5F8A}
\definecolor{accentGreen}{HTML}{2A7A2A}
\definecolor{accentOrange}{HTML}{CC6600}

% ── Theme ────────────────────────────────────────────────────────────────────
\usetheme{default}
\setbeamercolor{structure}{fg=cmuRed}
\setbeamercolor{block title}{fg=white, bg=cmuRed}
\setbeamercolor{block body}{fg=black, bg=lightGray}
\setbeamercolor{block title alerted}{fg=white, bg=accentBlue}
\setbeamercolor{block body alerted}{fg=black, bg=lightGray}
\setbeamerfont{block title}{size=\large, series=\bfseries}
\setbeamerfont{block body}{size=\normalsize}
\setbeamertemplate{navigation symbols}{}

\newcommand{\tight}{\vspace{-0.5em}}
\newcommand{\sep}{\vspace{0.8em}}

% ─────────────────────────────────────────────────────────────────────────────
\begin{document}
\begin{frame}[t]{}

%% ── HEADER ──────────────────────────────────────────────────────────────────
\begin{columns}[t]
  \begin{column}{0.82\linewidth}
    {\color{cmuRed}\Huge\bfseries SAPS: Step-Aware Pruning Schedule for Diffusion LLMs}\\[0.4em]
    {\Large Vedanti Kshirsagar \quad Sai Ravi Teja Gangavarapu \quad Riya Elizabeth John}\\[0.2em]
    {\large 17-752 Machine Learning Systems \ $\cdot$\ Carnegie Mellon University}
  \end{column}
  \begin{column}{0.16\linewidth}
    \raggedleft
    {\color{cmuGray}\large Model: LLaDA-8B-Instruct\\
                    Task: GSM8K (1{,}319 examples)\\
                    Platform: A100-80\,GB (Modal)}
  \end{column}
\end{columns}

\vspace{1.2em}
\rule{\linewidth}{2pt}
\vspace{0.8em}

%% ── BODY: 3 COLUMNS ─────────────────────────────────────────────────────────
\begin{columns}[t,totalwidth=\linewidth]

%%%% COL 1 ───────────────────────────────────────────────────────────────────
\begin{column}{0.30\linewidth}

  \begin{block}{The Problem: Fixed Pruning Ignores Time}
    Diffusion LLMs (dLLMs) generate text by iteratively denoising a
    masked sequence over $T$ steps. Because attention is bidirectional,
    the \textbf{full KV cache is held across every step} — a persistent
    memory bottleneck.

    \sep
    \textbf{Sparse-dLLM} addresses this with dynamic eviction, but keeps a
    \emph{constant} fraction of tokens (e.g.\ keep\_ratio\,$= 0.5$) at
    every single step.

    \begin{center}
    \begin{tikzpicture}
      \foreach \i in {1,...,6} {
        \fill[accentBlue!60] (\i*1.2-1.2, 0) rectangle ++(0.9, 0.7);
      }
      \node[below] at (3.0, -0.1) {\small Sparse-dLLM: same budget every step};
    \end{tikzpicture}
    \end{center}

    \sep
    {\color{cmuRed}\textbf{This is structurally wrong.}}\\
    The model's context needs are \emph{not} uniform over time.
  \end{block}

  \vspace{0.8em}

  \begin{block}{The Insight: Denoising Has Two Regimes}
    \begin{center}
    \begin{tikzpicture}[scale=1.1]
      % early bar — tall
      \fill[accentGreen!70] (0,0) rectangle (1.8, 3.2);
      \node[above] at (0.9, 3.2) {\small\bfseries Early steps};
      \node[below, align=center] at (0.9, -0.1)
        {\small Global structure\\\small Keep \textbf{70\%}};

      % arrow
      \draw[->, thick, cmuGray] (2.2, 1.6) -- (3.2, 1.6);

      % late bar — short
      \fill[cmuRed!70] (3.5, 0) rectangle (5.3, 0.46);
      \node[above] at (4.4, 0.46) {\small\bfseries Late steps};
      \node[below, align=center] at (4.4, -0.1)
        {\small Local refinement\\\small Keep \textbf{10\%}};
    \end{tikzpicture}
    \end{center}

    \begin{itemize}
      \item \textbf{Early} ($t \approx 0$): Sequence mostly masked.
            Model forms global plans — topic, clause structure,
            cross-sentence links. \emph{Needs broad context.}
      \item \textbf{Late} ($t \approx T$): Most tokens decided.
            Remaining positions resolve local details.
            Distal context is redundant. \emph{Safe to prune hard.}
    \end{itemize}
  \end{block}

\end{column}

%%%% COL 2 ───────────────────────────────────────────────────────────────────
\begin{column}{0.36\linewidth}

  \begin{alertblock}{SAPS: The Schedule}
    Replace the fixed retention ratio $r$ with a \textbf{monotonically
    decreasing schedule} $r(t)$:

    \[
      r(t) \;=\; r_{\max} \cdot \!\left(\frac{r_{\min}}{r_{\max}}\right)^{\!u},
      \quad u = \frac{t}{T-1}
    \]

    Gives $r_{\max}$ at $t=0$ and $r_{\min}$ at $t=T-1$, smooth
    exponential decay in between.

    \sep
    \begin{center}
    \begin{tikzpicture}[scale=0.95]
      % axes
      \draw[->] (0,0) -- (8.5,0) node[right] {\small step $t$};
      \draw[->] (0,0) -- (0,4.5) node[above] {\small $r(t)$};

      % SAPS curve (exp decay from 0.7 to 0.1 over 8 points)
      \draw[accentBlue, very thick]
        (0,3.5) -- (1,2.89) -- (2,2.33) -- (3,1.85) -- (4,1.46)
        -- (5,1.14) -- (6,0.87) -- (7,0.62) -- (8,0.45);
      \node[right, accentBlue] at (8, 0.45) {\small SAPS};

      % Sparse-dLLM flat line at 0.5
      \draw[cmuGray, thick, dashed]
        (0,2.50) -- (8,2.50);
      \node[right, cmuGray] at (8, 2.50) {\small Sparse-dLLM};

      % r_max / r_min labels
      \draw[dotted, cmuGray] (0,3.5) -- (8,3.5);
      \node[left] at (0,3.5) {\small $r_{\max}{=}0.7$};
      \draw[dotted, cmuGray] (0,0.45) -- (8,0.45);
      \node[left] at (0,0.45) {\small $r_{\min}{=}0.1$};
    \end{tikzpicture}
    \end{center}

    \tight
    \textbf{Parameters used:} $r_{\max}=0.7$,\; $r_{\min}=0.1$,\;
    \texttt{decay=exp}
  \end{alertblock}

  \vspace{0.8em}

  \begin{block}{Implementation: 3 Lightweight Patches}
    SAPS requires \textbf{no model retraining}. It patches the
    Sparse-dLLM inference pipeline at workspace setup time:

    \vspace{0.4em}
    \begin{tabular}{@{}lp{0.72\linewidth}@{}}
      \textbf{Patch 1} & \texttt{CustomCache.filter\_cache}: call
                          \texttt{ratio\_controller.keep\_num(n)}
                          instead of fixed \texttt{keep\_ratio~$\times$~n} \\[0.4em]
      \textbf{Patch 2} & \texttt{generate()}: compute global step index,
                          call \texttt{set\_step(t,\,T)} before each
                          forward pass \\[0.4em]
      \textbf{Patch 3} & \texttt{LLaDACausalLM} wrapper: instantiate
                          \texttt{RatioController} from \texttt{saps\_config}
                          and pass it through \\
    \end{tabular}

    \vspace{0.4em}
    Overhead: \textbf{one floating-point op per denoising step}.
  \end{block}

\end{column}

%%%% COL 3 ───────────────────────────────────────────────────────────────────
\begin{column}{0.30\linewidth}

  \begin{block}{Results (LLaDA-8B, GSM8K)}

    \begin{center}
    \begin{tabular}{lccc}
      \toprule
      \textbf{Method} & \textbf{KV Cache} & \textbf{Jaccard} & \textbf{Accuracy} \\
      \midrule
      Vanilla LLaDA   & $\sim$0.145\,GiB & — & 78.17\% \\
      Sparse-dLLM     & 0.0724\,GiB      & 0.633 & 76.3\% \\
      \rowcolor{accentGreen!18}
      \textbf{SAPS (ours)} & \textbf{0.0496\,GiB} & 0.498 & \textbf{78.2\%} \\
      \midrule
      \rowcolor{accentGreen!10}
      \textit{vs.\ Sparse} & {\color{accentGreen}\textbf{$-$31.4\%}} & — & {\color{accentGreen}\textbf{$+$1.9pp}} \\
      \bottomrule
    \end{tabular}
    \end{center}

    \vspace{0.5em}
    {\small Profiling: 128-example dev set.\;
            Accuracy: 1{,}319-example full test set.\;
            All runs: steps=256, seed=2025.}

    \vspace{0.6em}
    \textbf{Code benchmarks (n=500 MBPP, n=164 HumanEval):}\\[0.3em]
    {\small Both pruning methods degrade $\sim$5\,pp from vanilla on code.
    SAPS and Sparse are statistically tied (MBPP: $\Delta{=}{-}0.8$\,pp,
    $p{=}0.78$). Step-aware scheduling helps \emph{math}, not code,
    at $r_{\min}{=}0.1$.}

    \vspace{0.8em}
    {\large\color{cmuRed}\textbf{Pareto improvement:}\\
    Less memory \emph{and} better accuracy — on math.}
  \end{block}

  \vspace{0.8em}

  \begin{block}{Why Does Accuracy Improve?}
    \begin{itemize}
      \item \textbf{Better early retention (70\%)} gives the model
            richer global context when forming the overall answer
            structure — leading to more coherent multi-step reasoning.

      \item \textbf{Aggressive late pruning (10\%)} forces attention
            onto the most salient nearby context, suppressing noise
            in the final token decisions (inference-time regularization).
    \end{itemize}

    \vspace{0.3em}
    SAPS \textbf{matches vanilla} LLaDA accuracy ($78.2\%$ vs $78.17\%$,
    $\Delta{=}{+}0.03$\,pp) while using $\sim$66\% less KV memory —
    recovering full unpruned quality for free.
  \end{block}

  \vspace{0.8em}

  \begin{block}{Takeaway}
    \textbf{The information a diffusion LLM needs from its KV cache is
    not constant over time.} A schedule that respects this temporal
    structure improves both efficiency and quality — no training required.

    \vspace{0.4em}
    \textit{Future work:} layer-aware schedule $r(t, \ell)$;
    learned schedules; longer-context evaluation.
  \end{block}

\end{column}

\end{columns}

\end{frame}
\end{document}
